% ja/book.tex -- the master file for the Japanese edition.
% Hand-edited. Build with: finishing/tools/build_edition.sh ja
%
% Shared by every edition: the preamble (switched by \booklanguage,
% \booklocale and \bookmainfont), edition-common.tex (English captions, notes,
% heads) and edition-body.tex (the document). ja/sections.tex and
% ja/status.tex are GENERATED by finishing/tools/editions.py.
% Japanese is set by luatexja in Noto Serif CJK JP (a system font); Latin keeps the English book's face.
% No babel locale is loaded (none for \booklocale): the CJK locales fight the
% typesetting package, so translated sections are set under English, which here
% governs only Latin hyphenation, and the heads and captions are the edition's own.
% biblatex has no Japanese strings; the edition's own are in ja/biblatex-ja.lbx, mapped
% onto English so that every citation uses them.
%
% The title, subtitle and notice are a language model's drafts, not yet
% reviewed by a native reader; the title itself is an open question (D-810).
\documentclass[11pt,oneside]{book}
\def\bookcode{ja}
\def\booklanguage{english}
\input{ja/status}
\newcommand{\booklegal}{%
     著作権 \textcopyright{} 2026 \bookauthor。\\
     CC BY-NC-ND 4.0 ライセンスのもとで公開。\\[0.4\baselineskip]
     本書に引用した第三者の資料（題辞を含む）は各権利者に帰属し、\\
     フェアユースの原則に基づいて転載している。上記のライセンスはこれらには及ばない。\\[0.4\baselineskip]
     日本語訳は作業中である。未翻訳の節は英語のまま掲載している。\\
     翻訳済みの節は言語モデルが作成した草稿であり、まだ母語話者の確認を経ていない。%
     \ifx\editionstalelist\empty\else\\[0.4\baselineskip]
     以下の節は翻訳を更新中のため英語で掲載している：\\
     \editionstalelist。\fi}
\input{preamble}
% luatexja's default faces are Japanese ones not installed here; name the
% Noto CJK faces before it loads, or it stops on a missing font.
\makeatletter
\def\ltj@stdmcfont{name:NotoSerifCJKjp-Regular}
\def\ltj@stdgtfont{name:NotoSansCJKjp-Regular}
\makeatother
\usepackage{luatexja-fontspec}
\setmainjfont{Noto Serif CJK JP}[JFM=ujis, Language=Japanese,
  BoldFont=Noto Serif CJK JP Bold, ItalicFont=Noto Serif CJK JP,
  BoldItalicFont=Noto Serif CJK JP Bold]
\addto\captionsenglish{\renewcommand{\contentsname}{目次}%
  \renewcommand{\appendixname}{付録}\renewcommand{\bibname}{参考文献}}
\newcommand{\edchapnum}{第\thechapter 章}
\newcommand{\edappnum}{付録\thechapter}
\newcommand{\edpartnum}{第\arabic{part}部}
\DeclareLanguageMapping{english}{ja/biblatex-ja}
\input{edition-common}

\newcommand{\booktitlemain}{アンチファシスト・インテリジェンス}
\newcommand{\booksubtitle}{拒否するかもしれない機械と、それを形づくるべき公衆}
\newcommand{\bookauthor}{Joshua G. Stern}
\newcommand{\bookreferences}{参考文献}
\title{\booktitlemain：\booksubtitle}
\author{\bookauthor}
\date{}
\hypersetup{pdftitle={\booktitlemain：\booksubtitle}, pdfauthor={\bookauthor}}
\input{edition-body}
